% 02_related.tex  --  Related Work

\section{Related Work}
\label{sec:related}

\paragraph{Dataset distillation and condensation.}
Dataset distillation~\cite{wang2018datasetdistillation} and its successors
optimise a small synthetic set so that a model retrained on it matches one
trained on the full data, using gradient matching~\cite{zhao2021dcgm},
distribution matching~\cite{zhao2023dm}, or trajectory
matching~\cite{cazenavette2022mtt}. Comprehensive surveys~\cite{lei2023survey,
yu2023review,sachdeva2023survey} emphasize downstream evaluation, mostly on
image benchmarks, and do not provide a GBDT-specific split-landscape fidelity
metric. Extending distillation to tabular data has
attracted recent
attention~\cite{kang2024tabdistill,kang2025tabdistill,herurkar2024tabdist,
medvedev2020tabular}, but these works retain the retrain-and-evaluate paradigm.
Our contribution is orthogonal: rather than a new condensation method, we
introduce a no-candidate-retraining diagnostic that measures whether \emph{any}
condensed set, regardless of the method that produced it, preserves the GBDT
split-gain landscape.

\paragraph{Coreset selection.}
Classical coreset methods select real training examples by geometric or
gradient-based criteria: herding~\cite{welling2009herding} matches
distribution moments; k-center~\cite{sener2018kcenter} minimises the maximum
coverage distance; CRAIG~\cite{mirzasoleiman2020craig} and
GRAD-MATCH~\cite{killamsetty2021gradmatch} match gradient norms via submodular
maximisation; EL2N and GraNd~\cite{paul2021el2n} use early-training gradient
norms. All of these methods were designed for and evaluated on differentiable
models. Translating their selection signals to the GBDT setting is non-trivial
because GBDTs are not differentiable with respect to the training set; we
include adapted variants in our 17-method benchmark and show that their
landscape fidelity, measured by SLD, predicts their tabular AUROC. The
GBDT-native subsampling methods GOSS~\cite{ke2017lightgbm} and
MVS~\cite{ibragimov2019mvs} operate per boosting round rather than producing
a static condensed set; MVS is the theoretically closest prior work, as it
explicitly minimises the variance of the split-scoring estimator. Data
valuation~\cite{ghorbani2019shapley,koh2017influence} and influence-function
methods offer related but complementary perspectives on data contribution; they
often require retraining on many subsets or differentiable-model assumptions
and are not directly applicable in their standard form to histogram GBDTs.

\paragraph{Why GBDTs for tabular data.}
Gradient-boosted decision trees~\cite{friedman2001gbm,chen2016xgboost,
ke2017lightgbm,prokhorenkova2018catboost} remain one of the strongest practical
model classes for tabular data~\cite{grinsztajn2022why,borisov2022survey}. This
motivates a condensation evaluation framework specific to tree-structured
inductive bias. The split-gain criterion~\cite{loh2011cart,natekin2013tutorial}
aggregates per-bin gradient statistics, and the resulting landscape fully
determines which feature and threshold a node selects, a structure that has no
direct analogue in differentiable models. Our SLD exploits this structure to
build a measurement instrument that is GBDT-mechanism-specific: it explains
\emph{why} a condensed set preserves tree behaviour, not merely whether it does.

\paragraph{The C\textsuperscript{2}TC contrast.}
The most closely related recent work is C\textsuperscript{2}TC~\cite{xu2026c2tc}
(ICDE 2026), a training-free tabular condensation
\emph{method}: it optimises class allocation and feature representation via
class-adaptive clustering to produce a condensed set efficiently. Our work is
fundamentally different in intent. C\textsuperscript{2}TC is a
\emph{condenser}: it produces a surrogate dataset. SLD is a
\emph{measurement instrument}: it audits the fidelity of any surrogate,
including one produced by C\textsuperscript{2}TC, with respect to the GBDT
split-gain landscape. Where C\textsuperscript{2}TC is model-agnostic in its
clustering objective, SLD is GBDT-mechanism-specific: it reads the landscape
from histogram statistics under a fixed prediction state without retraining on
each candidate. We do not include C\textsuperscript{2}TC in the full grid
because its official implementation was not integrated before our 17-method
grid was frozen. As a compatibility audit, we ran the official
C\textsuperscript{2}TC code on our Adult split at budget 200; the generated set
obtained AUROC $0.769$ and raw $\mathrm{SLD}_\infty=1228.2$, confirming that
SLD can audit C\textsuperscript{2}TC-generated surrogates; this single run does
not constitute a full C\textsuperscript{2}TC baseline.
